\documentclass[letterpaper]{article}
\PassOptionsToPackage{unicode=true}{hyperref} % options for packages loaded elsewhere
\PassOptionsToPackage{hyphens}{url}
\usepackage{graphicx} % Required for inserting images


\def\tightlist{}


\graphicspath{ {../images/} }
\usepackage[margin=1in]{geometry}
\usepackage{tikz}
\usepackage[dvipsnames]{xcolor}
\usetikzlibrary{calc}




\usepackage[utf8]{inputenc}
\usepackage[T1]{fontenc}
\usepackage{lmodern}
\usepackage[english]{babel}


% Useful packages
\usepackage{hyperref}
\usepackage{amsmath, amssymb}
\usepackage{fancyhdr}
\usepackage{setspace}
\usepackage{tocloft}
\usepackage{titlesec}


\usepackage{unicode-math}



% Code highlighting (Pandoc uses listings or minted, but listings is safer without shell escape)
\usepackage{listings}
\lstset{
  basicstyle=\ttfamily\small,
  breaklines=true,
  frame=single,
  backgroundcolor=\color[gray]{0.95}
}



%Modifiable Commands; defaults set to none
\newcommand{\titleDOC}{}
\newcommand{\authorDOC}{Daniel Sabalakov}
\newcommand{\dateDOC}{\today}
\newcommand{\classDOC}{}
% DEFAULT QUOTE IS BY ISSAC NEWTON. IF YOU DON'T WANT IT, DO: quote: @none
\newcommand{\quoteDOC}{``\textit{Truth is ever to be found in the simplicity, and not in the multiplicity and confusion of things}'' - Issac Newton}

% Header/Footer
\pagestyle{fancy}
\fancyhf{}
\rhead{\thepage}
\lhead{\leftmark}





% %Title Arguments
% \title{\TITLEOFDOCUMENT}
% \author{\AUTHOROFDOCUMENT}
% \date{\today}


%
% ANYTHING BELOW HERE CAN GET PROCEDURALLY GENERATED
%
% BEGINNING OF DOCUMENT
%
%
%
\begin{document}

%titlepage








\renewcommand{\titleDOC}{Down the River}

\renewcommand{\authorDOC}{Daniel Sabalakov, Anthony Smith, Soren Westerberg}

\renewcommand{\classDOC}{DE Intro to Statistics I}





\begin{titlepage}
  \titleDOC
  \classDOC
  \dateDOC
  \authorDOC
\end{titlepage}

Instructions: Use the graphical displays and statistical summaries to answer the questions. Recall, the farmer is going to use the results of a sample (n = 10) to determine whether it is worthwhile to harvest all plots (N = 100).

\begin{figure}[htbp]
	 \centering
	 \includegraphics[width=0.9\textwidth]{C:/Users/score/Downloads/s1.png}
\end{figure}

\begin{figure}[htbp]
	 \centering
	 \includegraphics[width=0.9\textwidth]{C:/Users/score/Downloads/s2.png}
\end{figure}

\subsubsection{Use the summary statistics to make a compelling case for why one design is better than the other two when comparing the SRS, Stratified by Column, and Stratified by Row sample designs.}\label{use-the-summary-statistics-to-make-a-compelling-case-for-why-one-design-is-better-than-the-other-two-when-comparing-the-srs-stratified-by-column-and-stratified-by-row-sample-designs.}

Stratified by Column is the best design to use for the problem Rolling Down the River. Stratified by Column has a mean of 220.60, and a standard deviation of 6.79, while Stratified by Row's mean is 223.82, and its standard deviation is 13.3. Lastly, the Simple Random Sample yields a mean of 220.68, and a standard deviation of 17.29. While all three designs have similar means, the difference is all in the spread. While Stratified by Column has a lot of outliers, this is only due to the outlier fence rule; almost all of Column's outliers lie within the 75\% of stratified by Row. Furthermore, since the IQR is low, 50\% of the data in Stratified by Column lies in a very small spread (218-221), while SRS and Row lie in (208.5-233.5) and (219-233) respectively. By stratifying by column, each level gets equal representation.



\end{document}
